\begin{afterword}
\summaryhead

The post sets out to ``restore a naive view of truth.'' When someone answers a clinical study by asking what ``true'' means, the author says, the question need not stop the argument, just as people fell off cliffs before anyone knew the equations of gravity. The rest is a dialogue.

A shepherd who keeps losing sheep to wolves builds a system: a pebble goes into a bucket for each sheep that leaves the fold and comes out for each that returns, so an empty bucket means an empty pasture. Mark, an official, wants to know where the ``magic'' is. He fills his own bucket to a level he likes, calls magic subjective, proposes that it comes from buckets agreeing with one another, and says sheep cannot affect buckets. The shepherd and the apprentice Autrey answer that the pebbles are ordinary, that what matters is a chain of cause and effect from sheep to pebble, and that adding pebbles does not create sheep. The system ``works whether or not you believe in it,'' because reality controls the belief and not the reverse. ``Reality'' is the name for whatever determines results that surprise even our best hypotheses.

Mark then argues that there is only belief, that truth is a tool of power, and that each culture has its own, and admits these are excuses he applies selectively. Inspector Darwin suggests Mark believe he can fly and step off a cliff. Mark falls. The shepherds go back to the sheep.

\responsehead

The dialogue gives a clear picture of how a belief comes to match the world. The count stays right because each sheep causes a pebble, and Mark's hand works as well as a bucket. What a pebble stands for is its place in the counting, not anything in the stone. The shepherd's account of ``reality,'' the name for whatever decides results that even good hypotheses fail to predict, is a fair answer to Mark's suggestion that ``true'' only means ``really really believe.'' Several of these ideas are known in philosophy: Peirce's 1878 definition of the real, teleological theories of content, and Stove's parody of arguments like Mark's. The post names none of them, though it does not claim them as new.

What the dialogue mainly argues for is realism: the world exists apart from our beliefs and causes them. Its thesis sentence says so: the correspondence ``comes from reality controlling my beliefs.'' That is not the same claim as the correspondence theory of truth, and some of that theory's rivals, deflationism among them, typically accept realism. The Tarski sentence the shepherd relies on is common ground among all of them. So the post defends its ``naive view of truth'' against people who deny reality, not against competing theories of truth.

The model covers one kind of case: a count of visible objects, checked the next morning. The standard first objection to correspondence theories is that they fit such cases and fit poorly on others, such as moral or logical truths. The post promises a view of truth without restriction and does not take up that objection. Arithmetic, its example of an obvious truth, is handled only as counting objects.

The opponent is a composite. Mark holds several views at once, and the dialogue itself notes that they are separate claims, then has Mark admit they are excuses. The main ones appear in the forms that the standard objections were made against. His coherentism is the crudest version: truth as buckets agreeing with one another. His relativism is the global kind, which the self-refutation argument applies to and which is much less common in current debates than local relativism. Where Mark can only name a problem, the philosophers whose labels he uses offer answers. As satire of arguers on mailing lists, which the post says is its aim, this is fair. As a reply to the positions whose names it uses, it does not engage them.

In short: a clear parable of how beliefs come to match the world, which defends realism more than the correspondence theory it names, on one kind of case, against an opponent written to hold the weakest forms of the views he names.
\end{afterword}
